\subsection{The exact remaining decision and comparisons}

\begin{statement}{AO.1}{Open}{The main quadratic decision on the full profinite pentagon group}
We have not proved or refuted
\[
                    \Delta|_{\GT}=0.
\]
With $\mathcal H_5=\GT\cap\HH$, the exact remaining identity
can now be written
\[
 \epsilon_2(W_h)=\kappa_d(h)\pmod2
                      \qquad(h\in\mathcal H_5),
\]
where $\kappa_d$ is the canonical integral character of FP.32.
Equivalently, by the genuine quadratic pentagon consequence
FP.37, the remaining identity is
\[
                 \nu_4(h)=\kappa_d(h)\pmod2
                    \qquad(h\in\mathcal H_5).
\]
Here $\nu_4$ is the explicit continuous $S_4$ evaluation
character, not a choice of a finite frame primitive.
No assertion in this return assigns a pentagon value to the R2
element, its cube, or the element of $\ker D$ in LI.3.
\end{statement}
The equivalence with the full-multiplier question is exact.
For $g\in\GT$, choose $\sigma\in G_{\Q}\subseteq\GT$ with
$\chi(\sigma)=\lambda(g)$. Then $h=g\sigma^{-1}\in\mathcal H_5$
and $\Delta(g)=\delta(h)$. Conversely a counterexample in
$\mathcal H_5$ is already one in $\GT$. FP.32 proves
$\delta(h)=\kappa_d(h)+\epsilon_2(W_h)$, giving the displayed
form. On the Galois subgroup both bits equal the quadratic
character of $\Q^{\rm ab}(\sqrt[4]3)/\Q^{\rm ab}$; on every
closed AA subgroup vanishing follows from A.3-L1.

The progress and the exact failures of the tested routes are as
follows. The pointed category kills every commutator parameter,
so its coherence cannot distinguish the R2 witness. The explicit
finite characteristic quotient of FP.23--FP.29 does retain the
actual frame signs and the outer-action ambiguity. However,
FP.30 proves that quotients induced only from forgetting points
annihilate every two-cycle pentagon defect. The genuine affine
Brunnian quotient of FP.31 imposes the proved equation
$a_f=\rho\pmod3$, but the actual R2 cube satisfies that equation
at modulus twelve while retaining $\Delta=1$ (FP.33).
Thus that equation does not force the displayed quadratic equality.
FP.37 now adds the actual quadratic consequence
$\epsilon_2(W)=\nu_4$; its finite cubic-frame ambiguity explains
exactly why this does not yet identify $\nu_4$ with $\kappa_d$.
The next required comparison must retain additional five-point
information, or identify these last two characters by an
independent theorem.

The mixed beta reduction MD.4 expresses the same problem as
$\Delta=d_2+e_3$, so the required conclusion is $d_2=e_3$;
proving both zero is sufficient but is a stronger assertion.
MD.8 proves that either descent fixed subgroup kills $d_2$.
MD.10 supplies an actual gentle element with every cyclic-cover
Kummer coefficient zero, beta series one, and nevertheless
$\Delta=1$. Thus even the entire coefficient package leaves
the cubic frame comparison essential.

LI.8 gives a single actual $k\in\mathcal H$ with
$\Delta(k)=1$, $D(k)=1$, and
$f_k\in R_{\rm nil}(F'')\subseteq F'''$.
FP.36 and FP.38 prove that it passes all pentagon evaluations
in the specified pure-solvable-length-three quotients, the
entire two-adic Burau family, and the first relative
characteristic-three homology extension of the order-$216$
five-point quotient. These are exact failures of those tests,
not hypothetical finite survivors. A new test must escape
their nilpotent second-derived evaluation images.
LI.9 strengthens the Burau conclusion to the entire adelic
congruence family at a primitive cube root of unity, using
one actual witness and exact integral identities.
FR.6 strengthens that construction simultaneously to the
whole quotient $L/R_{\rm nil}(C')$, where
$C=\ker(L\to\widehat B)=C_2\amalg C_2\amalg C_2$.
It retains the full profinite forgetful base. FR.7 also proves
that the corresponding braid formulas define an actual
automorphism of this quotient, fixing its Ferrari kernel
pointwise. Thus every finite five-point test with nilpotent
image of $C'$ is insufficient for this implication.
The next cover uses $D=C''(C')^2$, free of rank $129$, and
the actual action on $H_1(D,\mathbb F_3)$. FP.45 finds a
nonnilpotent harmonic second-derived image here, independently
checked by a second program. Nevertheless RP.8 proves
$\mathcal P(f)\in D$ for every actual multiplier-one gentle
element. Its conjugation action on every $H_1(D,R)$ is
therefore identity. Thus even this nonnilpotent linear test
is universally insufficient. The next precise target is
$L/(D'D^3)$, where the residual class in $D/(D'D^3)$
survives as an affine translation. Linear homology discards
exactly that translation; retaining it is a distinct test.
RA.3 now constructs an actual mixed-sign gentle element
passing this affine test as well, simultaneously with the
earlier complete families. Its cyclic seed has nilpotent
normal image in every evaluation and every peripheral core;
this is a proved input to R2, not a finite-shadow lift.
RA.4 upgrades it to an actual automorphism of
$L/(D'D^3)$ fixing $C/(D'D^3)$ pointwise.
RA.5 supplies that separate higher-level argument: exact
transfer isolates the rank-five factor, and congruence
kernels control its complementary affine factor. One
actual witness therefore satisfies the pentagon in the
whole quotient $L/D'$, and RA.6 proves its actual
automorphism there, fixing $C/D'$ pointwise. This retains
every coefficient level of $D^{\rm ab}$; it does not
retain the nonabelian pro-prime completions of $D$.
The next specified cover is $D_4=C''(C')^4$. Its full
Fox chain representations retain the augmentation line
as well as the homology, so their translation information
is not silently discarded.
The canonical homology detector RH.2 and its full-multiplier
bound LH.1 do detect $\Delta$ as an action character, but no
pentagon vanishing theorem for that action image is proved.

The mixed-sign restricted lifting problem in LX.2 is still
unproved: R1 and R2 do not lift the five correlated evaluations
of one pentagon parameter. We have no compatible system giving
an actual mixed-sign element of $\GT$. Nor have we computed a
finite detector at which the pentagon excludes every actual
gentle residue of sign one. The smallest such level is unknown.
These are the precise two missing alternatives, not an inference
from the survival of any finite shadow.

If vanishing is eventually proved, PC.1 has already established
that $\GT\subsetneq\ker\Delta\subsetneq\GZ$.
If it fails, the restriction is onto $\Ftwo$, and its kernel is
an index-two closed subgroup of $\GT$ containing $G_{\Q}$.
The latter alternative would prove $G_{\Q}\ne\GT$; no element
establishing that alternative is claimed here.

\begin{statement}{AO.2}{Open}{The corrected compatibility subgroup and the first iterated character}
Put $C_L=\operatorname{Fix}(\operatorname{Ad}_c\circ D_P)$.
The assertions $\GT\subseteq C_L$, $C_L\subseteq\ker\Delta$,
and the proposed characterization of $C_L$ or $\ker\Delta$
by the harmonic and equianharmonic equations alone remain
undecided here. The literal-current assertions
$\GT\subseteq\operatorname{Fix}(D_P)$ and
$\operatorname{Fix}(D_P)=\ker\Delta$ are already false.
\end{statement}
A4.M1--A4.M6 prove the marking conversion and the exact
GF0/HF0 primitive comparisons. LI.3 proves $C_L$ proper,
without assuming any parity theorem for its elements.
The equality $C_L=\ker\Delta$ is also false: the element
$k^2$ from LI.3 is nonidentity by R7, has $\Delta(k^2)=0$,
and has descendant one, so lies outside either fixed subgroup.
The character $\chi_1=\Delta\circ D_P$ vanishes on Galois
and differs from $\Delta$. It is either zero or independent
of $\Delta$, by LI.4. If it is zero, then both fixed subgroups
lie in $\ker\Delta$. We have proved neither its zero value
nor a nonzero value on any actual gentle element.
The common-word finite calculation LC.5 retains both bits
but supplies neither missing conclusion. The stronger
$D_P(k)=1$, $\Delta(k)=1$ theorem proves noninjectivity of
descent, not vanishing of the character on its image.
MD.8 does prove that both fixed subgroups annihilate the mixed
duplication defect $d_2$. Their remaining $\Delta$ comparison
is precisely the cubic defect, not the already settled
harmonic half-Kummer identity.

\begin{statement}{AO.3}{Open}{The well-defined level-six Atkin--Lehner comparison}
For the crossed defect $AL_{6,3}$ defined in FP.17, the
equivalence $AL_{6,3}(g)=1\Longleftrightarrow\Delta(g)=0$
on its specified gentle domain has not been proved or refuted.
The actual free-group value of this defect on the R2 element
has not been computed. The analogous full-GT compatibility
of the two actions at the section with coordinate nine
also remains undecided.
\end{statement}
The rational four-cusp model and all three Atkin--Lehner maps
are explicit in FP.10. Arithmetic functoriality proves the
Galois comparison. FP.18 proves that the defect always fixes
the four cusps, so the proposed cusp permutation test cannot
detect its value. A calculation on a genuinely marked free
basis, with the section and conjugating paths retained,
is still required. The general $\Delta_N$ construction of
Lead 8 was refuted in A.7--A.8; it cannot fill this gap.
